\section{Source audit and the adopted curve}

Endler and Gallas classify period-seven chiral doublets through
\begin{equation}\label{eq:c7}
 C_7(\sigma)=\sigma^2-2\sigma-a.
\end{equation}
Their printed coordinate carrier is monic of degree seven in $x$.  All terms
except the final constant block are frozen exactly as printed.  The literal
last block is
\begin{equation}\label{eq:literal}
 -2a^3+6a^2+2a+3(a^3-4a^2+a-2)\sigma,
\end{equation}
whereas the block selected at the specialization of
Proposition~\ref{prop:correction} and generically certified in
Section~\ref{sec:neighbor} is
\begin{equation}\label{eq:corrected}
 -2a^3+6a^2+2a+3+(a^3-4a^2+a-2)\sigma.
\end{equation}

\begin{proposition}[Exact correction witness]\label{prop:correction}
Over $\F_{103}$, let $a=6$ and $\sigma=26$.  The adopted carrier has root
set
\[
 \{10,17,31,54,58,67,98\}.
\]
These values support two period-seven state cycles exchanged by the reversor
$R(x,y)=(y,x)$.  The literal printed carrier has root set $\{55,60\}$ and
therefore does not encode these cycles.
\end{proposition}

\begin{proof}
Apply $H_6(x,y)=(6-x^2-y,x)$ modulo 103 to
\begin{align*}
 &(10,54),(58,10),(31,58),(17,31),(98,17),(67,98),(54,67),\\
 &(10,58),(54,10),(67,54),(98,67),(17,98),(31,17),(58,31).
\end{align*}
Each row closes after seven steps; swapping the two entries maps one row to a
cyclic rotation of the other.  Both coordinate sums are 26, and
$26^2-2\cdot26=6$ in $\F_{103}$.  Direct Horner evaluation gives the two root
sets stated above.
\end{proof}

The literal carrier has $x$-discriminant of degree 42 after substituting
Eq.~\eqref{eq:c7}; in particular it does not have the factorization proved in
Section~\ref{sec:genus}.  Proposition~\ref{prop:correction} is sufficient to
reject the literal equation for the stated orbit.  The remainder of the paper
is a theorem about the explicitly frozen adopted polynomial, not about the
literal typeset equation.  The generic H\'enon carrier property is proved
constructively in Section~\ref{sec:neighbor}; equality with every component
of the full saturated exact-period-seven scheme is not asserted.

Put $a=\sigma^2-2\sigma$ in the adopted carrier and write the resulting
polynomial as $P(\sigma,x)$; its full expression appears in
Appendix~\ref{app:poly}.  Let $\Ccurve$ be the smooth projective normalization
of the affine plane curve $P=0$.  The base $C_7=0$ is itself the rational
$\sigma$-line.  The genus calculation concerns the seven-sheeted coordinate
cover $\Ccurve\to\Pone_\sigma$, not the chiral parameter line.

At $a=0$, equivalently $\sigma=0$ or 2, the conjugacy between
Eqs.~\eqref{eq:paper5} and~\eqref{eq:hamiltonian} degenerates.  The points of
$\Ccurve$ above those parameters belong to the algebraic closure of the
Hamiltonian model; we do not assign them a direct Paper-5 phase-space
interpretation.
